\chapter{OSPFv2, Single Area}\label{ch:ospfv2}
\bp{3.3}

\begin{outcomes}
By the end of this chapter you will be able to:
\begin{tight}
  \item \textbf{Explain} how OSPF forms adjacencies and why they sometimes do not.
  \item \textbf{Configure} single-area OSPFv2 with an explicit router ID, on both
        broadcast and point-to-point links.
  \item \textbf{Predict} and \textbf{influence} DR and BDR election.
  \item \textbf{Troubleshoot} the five mismatches that prevent an adjacency, using
        \cmd{show ip ospf neighbor} and \cmd{show ip ospf interface}.
\end{tight}
\end{outcomes}

\begin{prereq}
\chref{routingtable} for administrative distance and metrics; \chref{static} for
the contrast with a route you configure by hand.
\end{prereq}

\begin{keyterms}
OSPF $\bullet$ link-state $\bullet$ adjacency $\bullet$ neighbour $\bullet$
hello $\bullet$ router ID $\bullet$ area 0 $\bullet$ DR $\bullet$ BDR $\bullet$
DROTHER $\bullet$ cost $\bullet$ reference bandwidth $\bullet$ wildcard mask
$\bullet$ passive interface $\bullet$ FULL state $\bullet$ 2-WAY state
\end{keyterms}

\begin{examnote}[What the blueprint asks for, precisely]
Topic \tp{3.3} is ``configure single area OSPFv2 for IPv4 and OSPFv3 for IPv6''
with four sub-points: neighbour adjacencies \emph{excluding authentication},
point-to-point, broadcast with DR/BDR selection, and router ID. Multi-area OSPF,
authentication, virtual links and route summarisation are all out of scope. Learn
these four properly rather than the whole protocol thinly.
\end{examnote}

\section{What OSPF does that a static route cannot}

\begin{conceptbox}[Link state, in three sentences]
Every OSPF router floods a description of its own links to every other router in
the area. Each router therefore ends up with an identical map of the whole area,
which it stores in a \term{link-state database}. Each then runs the shortest-path
algorithm against that map to work out its own best route to everywhere.

Because every router has the whole map, a failure anywhere is known everywhere
within seconds, and every router recalculates independently and consistently.
That is the property a static route can never have.
\end{conceptbox}

\ccnafig{%
\begin{tikzpicture}[font=\scriptsize]
  \tikzset{
    dv/.style={rounded corners=3pt, draw=#1, fill=#1!8, text=#1!45!black,
               line width=0.9pt, minimum width=1.45cm, minimum height=1.0cm,
               align=center, font=\tiny\bfseries},
    w/.style={line width=1pt, neutral!65},
    db/.style={draw=dom3!60, fill=dom3!6, rounded corners=2pt, line width=0.7pt,
               font=\tiny, align=center, inner sep=3pt, text=dom3!45!black}}

  \node[dv=dom3] (r1) at (0,0)   {\faIcon{route}\\R1};
  \node[dv=dom3] (r2) at (5.2,0) {\faIcon{route}\\R2};
  \node[dv=dom3] (r3) at (10.4,0){\faIcon{route}\\R3};
  \draw[w] (r1) -- (r2);  \draw[w] (r2) -- (r3);

  \node[font=\tiny\bfseries, text=dom2!45!black] at (2.6,0.34) {FULL};
  \node[font=\tiny\bfseries, text=dom2!45!black] at (7.8,0.34) {FULL};
  \foreach \x in {2.6, 7.8}
    \draw[{Stealth[length=1.8mm]}-{Stealth[length=1.8mm]}, dom2!70,
          line width=0.8pt] (\x-0.75,-0.30) -- (\x+0.75,-0.30);
  \node[font=\tiny, text=neutral!85!black] at (2.6,-0.62) {hellos};
  \node[font=\tiny, text=neutral!85!black] at (7.8,-0.62) {hellos};

  \foreach \x in {0, 5.2, 10.4}
    \node[db] at (\x,-1.75) {link-state database\\\textbf{identical on all three}};

  \foreach \x in {0, 5.2, 10.4}
    \draw[dom3!45, line width=0.7pt, dash pattern=on 2pt off 2pt]
      (\x,-0.55) -- (\x,-1.3);

  \node[font=\tiny, text=neutral!85!black, align=center] at (5.2,-2.75)
      {Each router floods a description of \emph{its own} links.
       Every router therefore holds the same map ---\\
       and then computes its \emph{own} best path across it. Nobody is told what
       their routes are.};
\end{tikzpicture}}%
{Link state in one picture. The database is shared and identical; the routing
table each router builds from it is its own.}{ospflsdb}

\section{Adjacencies, and the five things that must match}

Two routers become neighbours by exchanging hello packets on a shared link. The
hello carries several parameters, and \textbf{if any of them disagree, no
adjacency forms} --- silently.

\begin{longtable}{@{}L{4.0cm}L{11.2cm}@{}}
\toprule
\rowcolor{primary!15}
\textbf{Must match} & \textbf{What goes wrong if it does not} \\
\midrule
\endhead
\textbf{Area ID} & Different areas on the same link: no adjacency. In single-area
OSPF everything is area 0 \\
\rowcolor{lightbg}
\textbf{Hello and dead intervals} & Defaults are 10~s and 40~s on broadcast
links. Change one end only and the adjacency never forms \\
\textbf{Subnet and mask} & The two interfaces must be in the same subnet. A mask
typo is a common cause \\
\rowcolor{lightbg}
\textbf{Network type} & Broadcast on one end and point-to-point on the other
prevents adjacency or leaves it stuck \\
\textbf{Authentication} & Out of scope for this exam, but real: mismatched keys
prevent adjacency \\
\bottomrule
\end{longtable}

Two further conditions are worth adding because they cause the same symptom: the
\textbf{router IDs must be unique}, and the interface must not be
\cmd{passive}.

\begin{definitionbox}[Router ID, and how it is chosen]
The \term{router ID} is a 32-bit value written like an IP address that uniquely
identifies the router in OSPF. It is selected in this order:
\begin{tightnum}
  \item The value configured with \cmd{router-id} --- always use this.
  \item Otherwise, the highest IP address on any \emph{loopback} interface.
  \item Otherwise, the highest IP address on any active physical interface.
\end{tightnum}
Left to itself the ID can change when an interface goes down, which forces the
adjacency to rebuild. Configure it explicitly and the problem disappears.
\end{definitionbox}

\section{Configuring single-area OSPFv2}

\begin{config}{OSPF on R1, area 0}
! A loopback gives a stable ID and a pingable address that never
! goes down with a physical interface.
interface Loopback0
 ip address 1.1.1.1 255.255.255.255
!
router ospf 10
 ! The process number is LOCAL. R2 may use a different one and
 ! the two will still peer -- unlike EIGRP or BGP.
 router-id 1.1.1.1
 !
 ! The network statement uses a WILDCARD mask -- the inverse of a
 ! subnet mask. It means "enable OSPF on any interface whose
 ! address falls in this range, and advertise that subnet".
 network 10.255.0.0 0.0.0.3 area 0
 network 192.168.10.0 0.0.0.255 area 0
 !
 ! Do not send hellos out of a LAN with no other routers on it.
 ! It saves bandwidth and removes an attack surface, while the
 ! subnet is still advertised.
 passive-interface Vlan10
 !
 ! Make cost reflect modern link speeds -- see the note below.
 auto-cost reference-bandwidth 10000
\end{config}

\begin{conceptbox}[Wildcard masks, quickly]
A wildcard mask is the subnet mask with every bit inverted: \cmd{0} means ``this
bit must match'', \cmd{255} means ``do not care''.

\begin{tight}
  \item \cmd{/24} $\rightarrow$ mask \cmd{255.255.255.0} $\rightarrow$ wildcard
        \cmd{0.0.0.255}
  \item \cmd{/30} $\rightarrow$ mask \cmd{255.255.255.252} $\rightarrow$ wildcard
        \cmd{0.0.0.3}
  \item A single interface $\rightarrow$ wildcard \cmd{0.0.0.0}
\end{tight}
The arithmetic is $255$ minus each octet of the subnet mask. Wildcards return in
\chref{acl}.
\end{conceptbox}

\begin{alertbox}[Reference bandwidth: change it everywhere or nowhere]
OSPF cost is $\text{reference bandwidth} \div \text{interface bandwidth}$, and
the default reference is 100~Mbps. That means a 1~Gbps link and a 10~Gbps link
both come out at cost 1 --- OSPF cannot tell them apart, and will happily choose
the slower one.

Raising the reference to 10000 (10~Gbps) restores the distinction. But cost is
compared across the area, so \textbf{every router must use the same reference
bandwidth} or the path calculations disagree and traffic takes strange routes.
Change it on all of them, or on none.
\end{alertbox}

\section{DR and BDR on broadcast links}

\begin{conceptbox}[Why a designated router exists]
On a broadcast segment with $n$ routers, full adjacencies between every pair
would be $n(n-1)/2$ relationships --- ten routers would need forty-five. Instead
one router is elected \term{DR} and one \term{BDR}, and everyone else becomes a
\term{DROTHER} adjacent only to those two. Flooding goes through the DR.

The election: \textbf{highest OSPF interface priority} wins; if priorities tie,
\textbf{highest router ID} wins. Default priority is 1. Priority \cmd{0} means
``never be DR''.

Two things surprise people. The election is \textbf{per segment}, not per router
--- a router can be DR on one link and DROTHER on another. And it is
\textbf{non-preemptive}: a new router with a better priority does not take over
until the current DR fails. To force it, clear the process or bounce the
interface.
\end{conceptbox}

\begin{config}{influencing the election, and avoiding it}
! Make this router the DR on the segment.
interface GigabitEthernet0/0
 ip ospf priority 255
!
! Never allow this one to be DR or BDR.
interface GigabitEthernet0/1
 ip ospf priority 0
!
! On a link with exactly two routers, skip the election entirely.
! Point-to-point forms an adjacency faster and elects nobody.
interface GigabitEthernet0/2
 ip ospf network point-to-point
\end{config}

\section{Verifying}

\begin{verify}{R1}{show ip ospf neighbor}
Neighbor ID     Pri   State           Dead Time   Address         Interface
2.2.2.2           1   FULL/DR         00:00:33    10.255.0.2      GigabitEthernet0/1
3.3.3.3           1   FULL/BDR        00:00:31    10.255.0.3      GigabitEthernet0/1
4.4.4.4           0   FULL/  -        00:00:38    10.255.1.2      GigabitEthernet0/2
\end{verify}

\begin{longtable}{@{}L{3.4cm}L{11.8cm}@{}}
\toprule
\rowcolor{primary!15}
\textbf{State} & \textbf{Meaning} \\
\midrule
\endhead
\cmd{FULL} & Databases synchronised. This is the healthy state \\
\rowcolor{lightbg}
\cmd{2-WAY} & Normal \emph{between two DROTHERs} on a broadcast segment --- they
see each other but do not need full adjacency. Not a fault \\
\cmd{EXSTART}, \cmd{EXCHANGE} & Mid-negotiation. Stuck here usually means an MTU
mismatch \\
\rowcolor{lightbg}
\cmd{INIT} & Hellos received but this router is not listed in them --- often
one-way communication or an access list \\
\cmd{FULL/ -} & Full, on a point-to-point link where there is no DR \\
\bottomrule
\end{longtable}

\begin{verify}{R1}{show ip ospf interface GigabitEthernet0/1}
GigabitEthernet0/1 is up, line protocol is up
  Internet Address 10.255.0.1/30, Area 0, Attached via Network Statement
  Process ID 10, Router ID 1.1.1.1, Network Type BROADCAST, Cost: 10
  Designated Router (ID) 2.2.2.2, Interface address 10.255.0.2
  Backup Designated router (ID) 3.3.3.3, Interface address 10.255.0.3
  Timer intervals configured, Hello 10, Dead 40, Wait 40, Retransmit 5
  Neighbor Count is 2, Adjacent neighbor count is 2
\end{verify}

This one output carries four of the five things that must match: area, network
type, hello and dead intervals, and the subnet. When an adjacency will not form,
run it on \emph{both} routers and compare line by line.

\begin{pitfall}
\begin{tight}
  \item \textbf{Assuming the process number must match.} It is locally
        significant. \cmd{router ospf 10} peers happily with \cmd{router ospf 65}.
  \item \textbf{Using a subnet mask in the network statement.} It takes a
        wildcard. \cmd{network 10.255.0.0 255.255.255.252} does not do what you
        expect.
  \item \textbf{Leaving the router ID to chance.} It changes when interfaces
        change, and every adjacency rebuilds.
  \item \textbf{Panicking at \cmd{2-WAY}.} Between two DROTHERs it is correct.
  \item \textbf{Changing reference bandwidth on one router.} See the warning
        above.
  \item \textbf{Forgetting \cmd{passive-interface} on user VLANs.} Hellos on a
        LAN full of PCs are wasted and expose the protocol.
\end{tight}
\end{pitfall}

\begin{breakfix}[Adjacency stuck in EXSTART.]
Two routers see each other but never reach FULL:

\begin{verify}{R1}{show ip ospf neighbor}
Neighbor ID     Pri   State           Dead Time   Address         Interface
2.2.2.2           1   EXSTART/DR      00:00:35    10.255.0.2      GigabitEthernet0/1
\end{verify}

Hellos are being exchanged --- so area, timers, subnet and network type all match,
or the neighbour would not be listed at all. The failure is later, during database
exchange.

\begin{verify}{R1}{show ip ospf interface GigabitEthernet0/1 | include MTU|Internet}
  Internet Address 10.255.0.1/30, Area 0, Attached via Network Statement
  MTU is 1500 bytes
\end{verify}

\begin{verify}{R2}{show ip ospf interface GigabitEthernet0/1 | include MTU|Internet}
  Internet Address 10.255.0.2/30, Area 0, Attached via Network Statement
  MTU is 1400 bytes
\end{verify}

\textbf{Diagnosis.} An MTU mismatch. OSPF checks MTU during database description
exchange and refuses to proceed if the two ends disagree --- which is exactly why
the adjacency reaches EXSTART and stops there. Being stuck in EXSTART or EXCHANGE
is close to diagnostic of this.

\textbf{Fix.} Make the MTUs match: \cmd{ip mtu 1500} on \cmd{R2}'s interface, or
1400 on \cmd{R1}'s if the smaller value was deliberate. The adjacency will move to
FULL within a dead interval. Do not use \cmd{ip ospf mtu-ignore} except as a
temporary measure --- it suppresses the symptom and leaves a real MTU mismatch to
fragment traffic later.
\end{breakfix}

\begin{lab}{Three routers, one area}{Packet Tracer}
\textbf{Topology.} \cmd{R1}, \cmd{R2}, \cmd{R3} on a shared Ethernet segment;
each with a LAN; \cmd{R1} also point-to-point to \cmd{R3}.

\begin{tasks}
  \item Configure loopbacks and explicit router IDs \cmd{1.1.1.1}, \cmd{2.2.2.2},
        \cmd{3.3.3.3}. Enable OSPF area 0 on all links and LANs, with LANs
        passive.
  \item Confirm every neighbour reaches FULL. Record which router became DR and
        which BDR, and explain the election from the router IDs.
  \item Set \cmd{ip ospf priority 255} on \cmd{R3}'s segment interface. Note
        that nothing changes, then bounce the interfaces and note that it does.
        Explain why.
  \item On the \cmd{R1}--\cmd{R3} link, set \cmd{ip ospf network point-to-point}
        at both ends. Confirm the state becomes \cmd{FULL/ -} with no DR.
  \item Change the hello interval on \cmd{R1} only. Watch the adjacency drop, and
        find the mismatch with \cmd{show ip ospf interface}.
  \item Set \cmd{ip mtu 1400} on one end and reproduce the EXSTART fault above.
  \item \textbf{Extension.} Raise the reference bandwidth on one router only and
        observe the asymmetric routing that results.
\end{tasks}
\end{lab}

\begin{checkpoint}
\begin{tightnum}
  \item Two routers on a broadcast segment show \cmd{2-WAY} and never go FULL.
        Under what circumstance is this correct rather than a fault?
  \item Which wildcard mask enables OSPF on \cmd{172.16.8.0/22}?
  \item A router's OSPF adjacencies all rebuild every time one interface flaps.
        What is the cause, and what is the one-line fix?
\end{tightnum}
\end{checkpoint}

\begin{chaptersummary}
\begin{tight}
  \item OSPF floods a description of every link, so all routers in the area share
        one map and converge consistently.
  \item Area, hello and dead timers, subnet and mask, and network type must all
        match, or no adjacency forms. Router IDs must be unique.
  \item Configure the router ID explicitly; otherwise it moves and adjacencies
        rebuild.
  \item Network statements take wildcard masks. The process number is local.
  \item DR and BDR are elected per segment by priority then router ID, and the
        election is non-preemptive. \cmd{ip ospf network point-to-point} avoids
        it entirely on two-router links.
  \item \cmd{FULL} is healthy; \cmd{2-WAY} between DROTHERs is healthy; stuck in
        \cmd{EXSTART} means MTU.
\end{tight}
\end{chaptersummary}

\begin{reviewq}
  \item \textbf{(MCQ)} Which is chosen as router ID if no \cmd{router-id} is
        configured and a loopback exists?
        (a)~lowest physical interface address (b)~highest loopback address
        (c)~highest physical address (d)~the process number
  \item \textbf{(MCQ)} Two routers are stuck in EXSTART. The most likely cause is
        (a)~area mismatch (b)~MTU mismatch (c)~timer mismatch (d)~duplicate router
        ID
  \item \textbf{(Short)} Explain why the OSPF process number does not need to
        match between neighbours, and name a protocol where the equivalent number
        does.
  \item \textbf{(Short)} Give the reason for setting \cmd{ip ospf priority 0} on
        an interface, and name a situation where you would.
  \item \textbf{(Applied)} Write a complete single-area OSPFv2 configuration for a
        router with a loopback, one point-to-point link to a neighbour and two
        user VLANs, and list the two \cmd{show} commands you would use to prove it
        is correct.
\end{reviewq}
